\Chapter{106}

\Verse{1}
{
    \zR{话说贾政闻知贾母危急，即忙进去看视。}
}
{
    \eR{[English source not present in the checked-in Bencraft/Joly files.]}
}
{
}

\Verse{2}
{
    \zR{见贾母惊吓气逆，王夫人鸳鸯等唤醒 回来，即用疏气安神的丸药服了，渐渐的好些，只是伤心落泪。}
}
{
    \eR{[English source not present in the checked-in Bencraft/Joly files.]}
}
{
}

\Verse{3}
{
    \zR{贾政在旁劝慰，总 说：“是儿子们不肖，招了祸来，累老太太受惊。}
}
{
    \eR{[English source not present in the checked-in Bencraft/Joly files.]}
}
{
}

\Verse{4}
{
    \zR{若老太太宽慰些，儿子们尚可在 外料理；若是老太太有什么不自在，儿子们的罪孽更重了。”}
}
{
    \eR{[English source not present in the checked-in Bencraft/Joly files.]}
}
{
}

\Verse{5}
{
    \zR{贾母道：“我活了八十 多岁，自作女孩儿起，到你父亲手里，都托着祖宗的福，从没有听见过这些事。}
}
{
    \eR{[English source not present in the checked-in Bencraft/Joly files.]}
}
{
}

\Verse{6}
{
    \zR{如 今到老了，见你们倘或受罪，叫我心里过的去吗？}
}
{
    \eR{[English source not present in the checked-in Bencraft/Joly files.]}
}
{
}

\Verse{7}
{
    \zR{倒不如合上眼随你们去罢了。”}
}
{
    \eR{[English source not present in the checked-in Bencraft/Joly files.]}
}
{
}

\Verse{8}
{
    \zR{说 着又哭。}
}
{
    \eR{[English source not present in the checked-in Bencraft/Joly files.]}
}
{
}

\Verse{9}
{
    \zR{贾政此时着急异常，又听外面说：“请老爷，内廷有信。”}
}
{
    \eR{[English source not present in the checked-in Bencraft/Joly files.]}
}
{
}

\Verse{10}
{
    \zR{贾政急忙出来，见是 北静王府长史，一见面便说：“大喜！”}
}
{
    \eR{[English source not present in the checked-in Bencraft/Joly files.]}
}
{
}

\Verse{11}
{
    \zR{贾政谢了，请长史坐下，请问：“王爷有何 谕旨？”}
}
{
    \eR{[English source not present in the checked-in Bencraft/Joly files.]}
}
{
}

\Verse{12}
{
    \zR{那长史道：“我们王爷同西平郡王进内复奏，将大人惧怕之心、感激天恩 之语都代奏过了。}
}
{
    \eR{[English source not present in the checked-in Bencraft/Joly files.]}
}
{
}

\Verse{13}
{
    \zR{主上甚是悯恤，并念及贵妃溘逝未久，不忍加罪，着加恩仍在工 部员外上行走。}
}
{
    \eR{[English source not present in the checked-in Bencraft/Joly files.]}
}
{
}

\Verse{14}
{
    \zR{所封家产，惟将贾赦的入官，馀俱给还，并传旨令尽心供职。}
}
{
    \eR{[English source not present in the checked-in Bencraft/Joly files.]}
}
{
}

\Verse{15}
{
    \zR{惟抄 出借券，令我们王爷查核。}
}
{
    \eR{[English source not present in the checked-in Bencraft/Joly files.]}
}
{
}

\Verse{16}
{
    \zR{如有违禁重利的，一概照例入官，其在定例生息的，同 房地文书，尽行给还。}
}
{
    \eR{[English source not present in the checked-in Bencraft/Joly files.]}
}
{
}

\Verse{17}
{
    \zR{贾琏着革去职衔，免罪释放。”}
}
{
    \eR{[English source not present in the checked-in Bencraft/Joly files.]}
}
{
}

\Verse{18}
{
    \zR{贾政听毕，即起身叩谢天恩， 又拜谢王爷恩典：“先请长史大人代为禀谢，明晨到阙谢恩，并到府里磕头。”}
}
{
    \eR{[English source not present in the checked-in Bencraft/Joly files.]}
}
{
}

\Verse{19}
{
    \zR{那长 史去了。}
}
{
    \eR{[English source not present in the checked-in Bencraft/Joly files.]}
}
{
}

\Verse{20}
{
    \zR{少停，传出旨来，承办官遵旨一一查清，入官者入官，给还者给还。}
}
{
    \eR{[English source not present in the checked-in Bencraft/Joly files.]}
}
{
}

\Verse{21}
{
    \zR{将贾 琏放出，所有贾赦名下男妇人等造册入官。}
}
{
    \eR{[English source not present in the checked-in Bencraft/Joly files.]}
}
{
}

\Verse{22}
{
    \zR{可怜贾琏屋内东西，除将按例放出的文书发给外，其馀虽未尽入官的，早被查 抄的人尽行抢去，所存者只有家伙物件。}
}
{
    \eR{[English source not present in the checked-in Bencraft/Joly files.]}
}
{
}

\Verse{23}
{
    \zR{贾琏始则惧罪，后蒙释放，已是大幸，及 想起历年积聚的东西并凤姐的体己，不下五七万金，一朝而尽，怎得不疼。}
}
{
    \eR{[English source not present in the checked-in Bencraft/Joly files.]}
}
{
}

\Verse{24}
{
    \zR{且他父 亲现禁在锦衣府，凤姐病在垂危，一时悲痛。}
}
{
    \eR{[English source not present in the checked-in Bencraft/Joly files.]}
}
{
}

\Verse{25}
{
    \zR{又见贾政含泪叫他，问道：“我因官 事在身，不大理家，故叫你们夫妇总理家事。}
}
{
    \eR{[English source not present in the checked-in Bencraft/Joly files.]}
}
{
}

\Verse{26}
{
    \zR{你父亲所为固难谏劝，那重利盘剥究 竟是谁干的？}
}
{
    \eR{[English source not present in the checked-in Bencraft/Joly files.]}
}
{
}

\Verse{27}
{
    \zR{况且非咱们这样人家所为。}
}
{
    \eR{[English source not present in the checked-in Bencraft/Joly files.]}
}
{
}

\Verse{28}
{
    \zR{如今入了官，在银钱呢是不打紧的，这声 名出去还了得吗！”}
}
{
    \eR{[English source not present in the checked-in Bencraft/Joly files.]}
}
{
}

\Verse{29}
{
    \zR{贾琏跪下说道：“侄儿办家事，并不敢存一点私心，所有出入的 帐目，自有赖大、吴新登、戴良等登记，老爷只管叫他们来查问。}
}
{
    \eR{[English source not present in the checked-in Bencraft/Joly files.]}
}
{
}

\Verse{30}
{
    \zR{现在这几年，库 内的银子出多入少，虽没贴补在内，已在各处做了好些空头，求老爷问太太就知道 了。}
}
{
    \eR{[English source not present in the checked-in Bencraft/Joly files.]}
}
{
}

\Verse{31}
{
    \zR{这些放出去的帐，连侄儿也不知道那里的银子，要问周瑞、旺儿才知道。”}
}
{
    \eR{[English source not present in the checked-in Bencraft/Joly files.]}
}
{
}

\Verse{32}
{
    \zR{贾 政道：“据你说来，连你自己屋里的事还不知道，那些家中上下的事更不知道了!我 这会子也不查问你。}
}
{
    \eR{[English source not present in the checked-in Bencraft/Joly files.]}
}
{
}

\Verse{33}
{
    \zR{现今你无事的人，你父亲的事和你珍大哥的事，还不快去打听 打听吗？”}
}
{
    \eR{[English source not present in the checked-in Bencraft/Joly files.]}
}
{
}

\Verse{34}
{
    \zR{贾琏一心委屈，含着眼泪，答应了出去。}
}
{
    \eR{[English source not present in the checked-in Bencraft/Joly files.]}
}
{
}

\Verse{35}
{
    \zR{贾政连连叹气，想道：“我祖父勤劳王事，立下功勋，得了两个世职，如今两 房犯事，都革去了。}
}
{
    \eR{[English source not present in the checked-in Bencraft/Joly files.]}
}
{
}

\Verse{36}
{
    \zR{我瞧这些子侄没一个长进的。}
}
{
    \eR{[English source not present in the checked-in Bencraft/Joly files.]}
}
{
}

\Verse{37}
{
    \zR{老天哪，老天哪!我贾家何至一 败如此!我虽蒙圣恩格外垂慈，给还家产，那两处食用自应归并一处，叫我一人那 里支撑的住？}
}
{
    \eR{[English source not present in the checked-in Bencraft/Joly files.]}
}
{
}

\Verse{38}
{
    \zR{方才琏儿所说，更加诧异，说不但库上无银，而且尚有亏空，这几年 竟是虚名在外。}
}
{
    \eR{[English source not present in the checked-in Bencraft/Joly files.]}
}
{
}

\Verse{39}
{
    \zR{只恨我自己为什么糊涂若此？}
}
{
    \eR{[English source not present in the checked-in Bencraft/Joly files.]}
}
{
}

\Verse{40}
{
    \zR{倘或我珠儿在世，尚有膀臂；宝玉虽 大，更是无用之物。”}
}
{
    \eR{[English source not present in the checked-in Bencraft/Joly files.]}
}
{
}

\Verse{41}
{
    \zR{想到那里，不觉泪满衣襟。}
}
{
    \eR{[English source not present in the checked-in Bencraft/Joly files.]}
}
{
}

\Verse{42}
{
    \zR{又想：“老太太若大年纪，儿子们 并没奉养一日，反累他老人家吓得死去活来，种种罪孽，叫我委之何人？”}
}
{
    \eR{[English source not present in the checked-in Bencraft/Joly files.]}
}
{
}

\Verse{43}
{
    \zR{正在独 自悲切，只见家人禀报：“各亲友进来看候。”}
}
{
    \eR{[English source not present in the checked-in Bencraft/Joly files.]}
}
{
}

\Verse{44}
{
    \zR{贾政一一道谢，说起：“家门不幸， 是我不能管教子侄，所以至此。”}
}
{
    \eR{[English source not present in the checked-in Bencraft/Joly files.]}
}
{
}

\Verse{45}
{
    \zR{有的说：“我久知令兄赦大老爷行事不妥，那边珍 爷更加骄纵。}
}
{
    \eR{[English source not present in the checked-in Bencraft/Joly files.]}
}
{
}

\Verse{46}
{
    \zR{若说因官事错误得个不是，于心无愧；如今自己闹出的，倒带累了二 老爷。”}
}
{
    \eR{[English source not present in the checked-in Bencraft/Joly files.]}
}
{
}

\Verse{47}
{
    \zR{有的说：“人家闹的也多，也没见御史参奏。}
}
{
    \eR{[English source not present in the checked-in Bencraft/Joly files.]}
}
{
}

\Verse{48}
{
    \zR{不是珍老大得罪朋友，何至如 此。”}
}
{
    \eR{[English source not present in the checked-in Bencraft/Joly files.]}
}
{
}

\Verse{49}
{
    \zR{有的说：“也不怪御史，我们听见说是府上的家人同几个泥腿在外头哄嚷出来 的。}
}
{
    \eR{[English source not present in the checked-in Bencraft/Joly files.]}
}
{
}

\Verse{50}
{
    \zR{御史恐参奏不实，所以诓了这里的人去，才说出来的。}
}
{
    \eR{[English source not present in the checked-in Bencraft/Joly files.]}
}
{
}

\Verse{51}
{
    \zR{我想府上待下人最宽的， 为什么还有这事？”}
}
{
    \eR{[English source not present in the checked-in Bencraft/Joly files.]}
}
{
}

\Verse{52}
{
    \zR{有的说：“大凡奴才们是一个养活不得的。}
}
{
    \eR{[English source not present in the checked-in Bencraft/Joly files.]}
}
{
}

\Verse{53}
{
    \zR{今儿在这里都是好 亲友，我才敢说。}
}
{
    \eR{[English source not present in the checked-in Bencraft/Joly files.]}
}
{
}

\Verse{54}
{
    \zR{就是尊驾在外任，我保不得。}
}
{
    \eR{[English source not present in the checked-in Bencraft/Joly files.]}
}
{
}

\Verse{55}
{
    \zR{你是不爱钱的，那外头的风声也不 好，都是奴才们闹的，你该提防些。}
}
{
    \eR{[English source not present in the checked-in Bencraft/Joly files.]}
}
{
}

\Verse{56}
{
    \zR{如今虽说没有动你的家，倘或再遇着主上疑心 起来，好些不便呢。”}
}
{
    \eR{[English source not present in the checked-in Bencraft/Joly files.]}
}
{
}

\Verse{57}
{
    \zR{贾政听说，心下着忙道：“众位听见我的风声怎样？”}
}
{
    \eR{[English source not present in the checked-in Bencraft/Joly files.]}
}
{
}

\Verse{58}
{
    \zR{众人道： “我们虽没见实据，只听得外头人说你在粮道任上，怎么叫门上家人要钱。”}
}
{
    \eR{[English source not present in the checked-in Bencraft/Joly files.]}
}
{
}

\Verse{59}
{
    \zR{贾政 听了，便说道：“我这是对天可表的，从不敢起这个念头。}
}
{
    \eR{[English source not present in the checked-in Bencraft/Joly files.]}
}
{
}

\Verse{60}
{
    \zR{只是奴才们在外头招摇 撞骗，闹出事来，我就耽不起。”}
}
{
    \eR{[English source not present in the checked-in Bencraft/Joly files.]}
}
{
}

\Verse{61}
{
    \zR{众人道：“如今怕也无益，只好将现在的管家们都 严严的查一查，若有抗主的奴才，查出来严严的办一办也罢了。”}
}
{
    \eR{[English source not present in the checked-in Bencraft/Joly files.]}
}
{
}

\Verse{62}
{
    \zR{贾政听了点头。}
}
{
    \eR{[English source not present in the checked-in Bencraft/Joly files.]}
}
{
}

\Verse{63}
{
    \zR{便见门上的进来回说：“孙姑爷打发人来说，自己有事不能来， 着人来瞧瞧。}
}
{
    \eR{[English source not present in the checked-in Bencraft/Joly files.]}
}
{
}

\Verse{64}
{
    \zR{说大老爷该他一项银子，要在二老爷身上还的。”}
}
{
    \eR{[English source not present in the checked-in Bencraft/Joly files.]}
}
{
}

\Verse{65}
{
    \zR{贾政心内忧闷，只 说：“知道了。”}
}
{
    \eR{[English source not present in the checked-in Bencraft/Joly files.]}
}
{
}

\Verse{66}
{
    \zR{众人都冷笑道：“人说令亲孙绍祖混帐，果然有的。}
}
{
    \eR{[English source not present in the checked-in Bencraft/Joly files.]}
}
{
}

\Verse{67}
{
    \zR{如今丈人抄了 家，不但不来瞧看帮补，倒赶忙的来要银子，真真不在理上。”}
}
{
    \eR{[English source not present in the checked-in Bencraft/Joly files.]}
}
{
}

\Verse{68}
{
    \zR{贾政道：“如今且不 必说他，那头亲事原是家兄配错了的。}
}
{
    \eR{[English source not present in the checked-in Bencraft/Joly files.]}
}
{
}

\Verse{69}
{
    \zR{我的侄女儿的罪已经受够了，如今又找上我 来了。”}
}
{
    \eR{[English source not present in the checked-in Bencraft/Joly files.]}
}
{
}

\Verse{70}
{
    \zR{正说着，只见薛蝌进来说道：“我打听锦衣府赵堂官必要照御史参的办，只 怕大老爷和珍大爷吃不住。”}
}
{
    \eR{[English source not present in the checked-in Bencraft/Joly files.]}
}
{
}

\Verse{71}
{
    \zR{众人都道：“二老爷，还是得你出去求求王爷，怎么挽 回挽回才好。}
}
{
    \eR{[English source not present in the checked-in Bencraft/Joly files.]}
}
{
}

\Verse{72}
{
    \zR{不然，这两家子就完了。”}
}
{
    \eR{[English source not present in the checked-in Bencraft/Joly files.]}
}
{
}

\Verse{73}
{
    \zR{贾政答应致谢，众人都散。}
}
{
    \eR{[English source not present in the checked-in Bencraft/Joly files.]}
}
{
}

\Verse{74}
{
    \zR{那时天已点灯时候，贾政进去请贾母的安，见贾母略略好些。}
}
{
    \eR{[English source not present in the checked-in Bencraft/Joly files.]}
}
{
}

\Verse{75}
{
    \zR{回到自己房中， 埋怨贾琏夫妇不知好歹，如今闹出放账的事情，大家不好，心里很不受用。}
}
{
    \eR{[English source not present in the checked-in Bencraft/Joly files.]}
}
{
}

\Verse{76}
{
    \zR{只是凤 姐现在病重，况他所有的什物尽被抄抢，心内自然难受，一时也未便说他，暂且隐 忍不言。}
}
{
    \eR{[English source not present in the checked-in Bencraft/Joly files.]}
}
{
}

\Verse{77}
{
    \zR{一夜无话。}
}
{
    \eR{[English source not present in the checked-in Bencraft/Joly files.]}
}
{
}

\Verse{78}
{
    \zR{次早贾政进内谢恩，并到北静王府西平王府两处叩谢，求二位 王爷照应他哥哥侄儿。}
}
{
    \eR{[English source not present in the checked-in Bencraft/Joly files.]}
}
{
}

\Verse{79}
{
    \zR{二王应许。}
}
{
    \eR{[English source not present in the checked-in Bencraft/Joly files.]}
}
{
}

\Verse{80}
{
    \zR{贾政又在同寅相好处托情。}
}
{
    \eR{[English source not present in the checked-in Bencraft/Joly files.]}
}
{
}

\Verse{81}
{
    \zR{且说贾琏打听得父兄之事不大妥，无法可施，只得回到家中。}
}
{
    \eR{[English source not present in the checked-in Bencraft/Joly files.]}
}
{
}

\Verse{82}
{
    \zR{平儿守着凤姐哭 泣，秋桐在耳房里抱怨凤姐。}
}
{
    \eR{[English source not present in the checked-in Bencraft/Joly files.]}
}
{
}

\Verse{83}
{
    \zR{贾琏走到旁边，见凤姐奄奄一息，就有多少怨言，一 时也说不出来。}
}
{
    \eR{[English source not present in the checked-in Bencraft/Joly files.]}
}
{
}

\Verse{84}
{
    \zR{平儿哭道：“如今已经这样，东西去了不能复来；奶奶这样，还得 再请个大夫瞧瞧才好啊。”}
}
{
    \eR{[English source not present in the checked-in Bencraft/Joly files.]}
}
{
}

\Verse{85}
{
    \zR{贾琏啐道：“呸!我的性命还不保，我还管他呢！”}
}
{
    \eR{[English source not present in the checked-in Bencraft/Joly files.]}
}
{
}

\Verse{86}
{
    \zR{凤姐听 见，睁眼一瞧，虽不言语，那眼泪直流。}
}
{
    \eR{[English source not present in the checked-in Bencraft/Joly files.]}
}
{
}

\Verse{87}
{
    \zR{看见贾琏出去了，便和平儿道：“你别不 达时务了。}
}
{
    \eR{[English source not present in the checked-in Bencraft/Joly files.]}
}
{
}

\Verse{88}
{
    \zR{到了这个田地，你还顾我做什么？}
}
{
    \eR{[English source not present in the checked-in Bencraft/Joly files.]}
}
{
}

\Verse{89}
{
    \zR{我巴不得今儿就死才好。}
}
{
    \eR{[English source not present in the checked-in Bencraft/Joly files.]}
}
{
}

\Verse{90}
{
    \zR{只要你能够 眼里有我，我死后你扶养大了巧姐儿，我在阴司里也感激你的情。”}
}
{
    \eR{[English source not present in the checked-in Bencraft/Joly files.]}
}
{
}

\Verse{91}
{
    \zR{平儿听了，越 发抽抽搭搭的哭起来了。}
}
{
    \eR{[English source not present in the checked-in Bencraft/Joly files.]}
}
{
}

\Verse{92}
{
    \zR{凤姐道：“你也不糊涂。}
}
{
    \eR{[English source not present in the checked-in Bencraft/Joly files.]}
}
{
}

\Verse{93}
{
    \zR{他们虽没有来说，必是抱怨我的。}
}
{
    \eR{[English source not present in the checked-in Bencraft/Joly files.]}
}
{
}

\Verse{94}
{
    \zR{虽说事是外头闹起，我不放账，也没我的事。}
}
{
    \eR{[English source not present in the checked-in Bencraft/Joly files.]}
}
{
}

\Verse{95}
{
    \zR{如今枉费心计，挣了一辈子的强，偏 偏儿的落在人后头了!我还恍惚听见珍大爷的事，说是强占良民妻子为妾，不从逼 死，有个姓张的在里头，你想想还有谁呢？}
}
{
    \eR{[English source not present in the checked-in Bencraft/Joly files.]}
}
{
}

\Verse{96}
{
    \zR{要是这件事审出来，咱们二爷是脱不了 的，我那时候儿可怎么见人呢？}
}
{
    \eR{[English source not present in the checked-in Bencraft/Joly files.]}
}
{
}

\Verse{97}
{
    \zR{我要立刻就死，又耽不起吞金服毒的。}
}
{
    \eR{[English source not present in the checked-in Bencraft/Joly files.]}
}
{
}

\Verse{98}
{
    \zR{你还要请大 夫，这不是你疼我，反倒害了我了么？”}
}
{
    \eR{[English source not present in the checked-in Bencraft/Joly files.]}
}
{
}

\Verse{99}
{
    \zR{平儿愈听愈惨，想来实在难处，恐凤姐自 寻短见，只得紧紧守着。}
}
{
    \eR{[English source not present in the checked-in Bencraft/Joly files.]}
}
{
}

\Verse{100}
{
    \zR{幸贾母不知底细，因近日身子好些，又见贾政无事，宝玉宝钗在旁，天天不离 左右，略觉放心。}
}
{
    \eR{[English source not present in the checked-in Bencraft/Joly files.]}
}
{
}

\Verse{101}
{
    \zR{素来最疼凤姐，便叫鸳鸯：“将我的体己东西拿些给凤丫头，再 拿些银钱交给平儿，好好的伏侍好了凤丫头，我再慢慢的分派。”}
}
{
    \eR{[English source not present in the checked-in Bencraft/Joly files.]}
}
{
}

\Verse{102}
{
    \zR{又命王夫人照看 邢夫人。}
}
{
    \eR{[English source not present in the checked-in Bencraft/Joly files.]}
}
{
}

\Verse{103}
{
    \zR{此时宁国府第入官，所有财产房地等项并家奴等俱已造册收尽。}
}
{
    \eR{[English source not present in the checked-in Bencraft/Joly files.]}
}
{
}

\Verse{104}
{
    \zR{这里贾母 命人将车接了尤氏婆媳过来。}
}
{
    \eR{[English source not present in the checked-in Bencraft/Joly files.]}
}
{
}

\Verse{105}
{
    \zR{可怜赫赫宁府，只剩得他们婆媳两个并佩凤偕鸾二人， 连一个下人没有。}
}
{
    \eR{[English source not present in the checked-in Bencraft/Joly files.]}
}
{
}

\Verse{106}
{
    \zR{贾母指出房子一所居住，就在惜春所住的间壁，又派了婆子四人、 丫头两个伏侍。}
}
{
    \eR{[English source not present in the checked-in Bencraft/Joly files.]}
}
{
}

\Verse{107}
{
    \zR{一应饭食起居在大厨房内分送，衣裙什物又是贾母送去，零星需用 亦在账房内开销，俱照荣府每人月例之数。}
}
{
    \eR{[English source not present in the checked-in Bencraft/Joly files.]}
}
{
}

\Verse{108}
{
    \zR{那贾赦、贾珍、贾蓉在锦衣府使用，账 房内实在无项可支。}
}
{
    \eR{[English source not present in the checked-in Bencraft/Joly files.]}
}
{
}

\Verse{109}
{
    \zR{如今凤姐儿一无所有，贾琏外头债务满身。}
}
{
    \eR{[English source not present in the checked-in Bencraft/Joly files.]}
}
{
}

\Verse{110}
{
    \zR{贾政不知家务，只 说：“已经托人，自有照应。”}
}
{
    \eR{[English source not present in the checked-in Bencraft/Joly files.]}
}
{
}

\Verse{111}
{
    \zR{贾琏无计可施，想到那亲戚里头，薛姨妈家已败，王 子腾已死，馀者亲戚虽有，俱是不能照应的，只得暗暗差人下屯，将地亩暂卖数千 金作为监中使费。}
}
{
    \eR{[English source not present in the checked-in Bencraft/Joly files.]}
}
{
}

\Verse{112}
{
    \zR{贾琏如此一行，那些家奴见主家势败，也便趁此弄鬼，并将东庄 租税也就指名借用些。}
}
{
    \eR{[English source not present in the checked-in Bencraft/Joly files.]}
}
{
}

\Verse{113}
{
    \zR{此是后话，暂且不提。}
}
{
    \eR{[English source not present in the checked-in Bencraft/Joly files.]}
}
{
}

\Verse{114}
{
    \zR{且说贾母见祖宗世职革去，现在子孙在监质审，邢夫人尤氏等日夜啼哭，凤姐 病在垂危，虽有宝玉宝钗在侧，只可解劝，不能分忧，所以日夜不宁，思前想后，
        眼泪不干。}
}
{
    \eR{[English source not present in the checked-in Bencraft/Joly files.]}
}
{
}

\Verse{115}
{
    \zR{一日傍晚，叫宝玉回去，自己扎挣坐起，叫鸳鸯等各处佛堂上香；又命 自己院内焚起斗香，用拐拄着，出到院中。}
}
{
    \eR{[English source not present in the checked-in Bencraft/Joly files.]}
}
{
}

\Verse{116}
{
    \zR{琥珀知是老太太拜佛，铺下大红猩毡拜 垫。}
}
{
    \eR{[English source not present in the checked-in Bencraft/Joly files.]}
}
{
}

\Verse{117}
{
    \zR{贾母上香跪下，磕了好些头，念了一回佛，含泪祝告天地道：“皇天菩萨在上： 我贾门史氏，虔诚祷告，求菩萨慈悲。}
}
{
    \eR{[English source not present in the checked-in Bencraft/Joly files.]}
}
{
}

\Verse{118}
{
    \zR{我贾门数世以来，不敢行凶霸道。}
}
{
    \eR{[English source not present in the checked-in Bencraft/Joly files.]}
}
{
}

\Verse{119}
{
    \zR{我帮夫助 子，虽不能为善，也不敢作恶。}
}
{
    \eR{[English source not present in the checked-in Bencraft/Joly files.]}
}
{
}

\Verse{120}
{
    \zR{必是后辈儿孙骄奢淫佚，暴殄天物，以致合府抄检。}
}
{
    \eR{[English source not present in the checked-in Bencraft/Joly files.]}
}
{
}

\Verse{121}
{
    \zR{现在儿孙监禁，自然凶多吉少，皆由我一人罪孽，不教儿孙，所以至此。}
}
{
    \eR{[English source not present in the checked-in Bencraft/Joly files.]}
}
{
}

\Verse{122}
{
    \zR{我今叩求 皇天保佑，在监的逢凶化吉，有病的早早安身。}
}
{
    \eR{[English source not present in the checked-in Bencraft/Joly files.]}
}
{
}

\Verse{123}
{
    \zR{总有合家罪孽，情愿一人承当，求 饶恕儿孙。}
}
{
    \eR{[English source not present in the checked-in Bencraft/Joly files.]}
}
{
}

\Verse{124}
{
    \zR{若皇天怜念我虔诚，早早赐我一死，宽免儿孙之罪！”}
}
{
    \eR{[English source not present in the checked-in Bencraft/Joly files.]}
}
{
}

\Verse{125}
{
    \zR{默默说到此处， 不禁伤心，呜呜咽咽的哭泣起来。}
}
{
    \eR{[English source not present in the checked-in Bencraft/Joly files.]}
}
{
}

\Verse{126}
{
    \zR{鸳鸯珍珠一面解劝，一面扶进房去。}
}
{
    \eR{[English source not present in the checked-in Bencraft/Joly files.]}
}
{
}

\Verse{127}
{
    \zR{只见王夫人带了宝玉宝钗过来请晚安，见贾母伤悲，三人也大哭起来。}
}
{
    \eR{[English source not present in the checked-in Bencraft/Joly files.]}
}
{
}

\Verse{128}
{
    \zR{宝钗更 有一层苦楚：想哥哥也在外监，将来要处决，不知可能减等；公婆虽然无事，眼见 家业萧条；宝玉依然疯傻，毫无志气。}
}
{
    \eR{[English source not present in the checked-in Bencraft/Joly files.]}
}
{
}

\Verse{129}
{
    \zR{想到后来终身，更比贾母王夫人哭的悲痛。}
}
{
    \eR{[English source not present in the checked-in Bencraft/Joly files.]}
}
{
}

\Verse{130}
{
    \zR{宝玉见宝钗如此，他也有一番悲戚，想着：“老太太年老不得安心，老爷太太见此 光景，不免悲伤，众姐妹风流云散，一日少似一日。}
}
{
    \eR{[English source not present in the checked-in Bencraft/Joly files.]}
}
{
}

\Verse{131}
{
    \zR{追思园中吟诗起社，何等热闹； 自林妹妹一死，我郁闷到今，又有宝姐姐伴着，不便时常哭泣。}
}
{
    \eR{[English source not present in the checked-in Bencraft/Joly files.]}
}
{
}

\Verse{132}
{
    \zR{况他又忧兄思母， 日夜难得笑容。}
}
{
    \eR{[English source not present in the checked-in Bencraft/Joly files.]}
}
{
}

\Verse{133}
{
    \zR{今日看他悲哀欲绝，心里更加不忍。”}
}
{
    \eR{[English source not present in the checked-in Bencraft/Joly files.]}
}
{
}

\Verse{134}
{
    \zR{竟嚎啕大哭起来。}
}
{
    \eR{[English source not present in the checked-in Bencraft/Joly files.]}
}
{
}

\Verse{135}
{
    \zR{鸳鸯、彩 云、莺儿、袭人看着，也各有所思，便都抽抽搭搭的。}
}
{
    \eR{[English source not present in the checked-in Bencraft/Joly files.]}
}
{
}

\Verse{136}
{
    \zR{馀者丫头们看的伤心，不觉 也都哭了。}
}
{
    \eR{[English source not present in the checked-in Bencraft/Joly files.]}
}
{
}

\Verse{137}
{
    \zR{竟无人劝。}
}
{
    \eR{[English source not present in the checked-in Bencraft/Joly files.]}
}
{
}

\Verse{138}
{
    \zR{满屋中哭声惊天动地，将外头上夜婆子吓慌，急报于贾政知 道。}
}
{
    \eR{[English source not present in the checked-in Bencraft/Joly files.]}
}
{
}

\Verse{139}
{
    \zR{那贾政正在书房纳闷，听见贾母的人来报，心中着忙，飞奔进内。}
}
{
    \eR{[English source not present in the checked-in Bencraft/Joly files.]}
}
{
}

\Verse{140}
{
    \zR{远远听得哭 声甚众，打量老太太不好，急的魂魄俱丧。}
}
{
    \eR{[English source not present in the checked-in Bencraft/Joly files.]}
}
{
}

\Verse{141}
{
    \zR{疾忙进来，只见坐着悲啼，才放下心来， 便道：“老太太伤心，你们该劝解才是啊，怎么打伙儿哭起来了？”}
}
{
    \eR{[English source not present in the checked-in Bencraft/Joly files.]}
}
{
}

\Verse{142}
{
    \zR{众人这才急忙 止哭，大家对面发怔。}
}
{
    \eR{[English source not present in the checked-in Bencraft/Joly files.]}
}
{
}

\Verse{143}
{
    \zR{贾政上前安慰了老太太，又说了众人几句。}
}
{
    \eR{[English source not present in the checked-in Bencraft/Joly files.]}
}
{
}

\Verse{144}
{
    \zR{都心里想道：“我 们原怕老太太悲伤，所以来劝解，怎么忘情，大家痛哭起来？”}
}
{
    \eR{[English source not present in the checked-in Bencraft/Joly files.]}
}
{
}

\Verse{145}
{
    \zR{正自不解，只见老婆子带了史侯家的两个女人进来，请了贾母的安，又向众人 请安毕，便说道：“我们家的老爷、太太、姑娘打发我来说：听见府里的事，原没
        什么大事，不过一时受惊。}
}
{
    \eR{[English source not present in the checked-in Bencraft/Joly files.]}
}
{
}

\Verse{146}
{
    \zR{恐怕老爷太太烦恼，叫我们过来告诉一声，说这里二老 爷是不怕的了。}
}
{
    \eR{[English source not present in the checked-in Bencraft/Joly files.]}
}
{
}

\Verse{147}
{
    \zR{我们姑娘本要自己来的，因不多几日就要出阁，所以不能来了。”}
}
{
    \eR{[English source not present in the checked-in Bencraft/Joly files.]}
}
{
}

\Verse{148}
{
    \zR{贾母听了，不便道谢，说：“你回去给我问好。}
}
{
    \eR{[English source not present in the checked-in Bencraft/Joly files.]}
}
{
}

\Verse{149}
{
    \zR{这是我们的家运合该如此。}
}
{
    \eR{[English source not present in the checked-in Bencraft/Joly files.]}
}
{
}

\Verse{150}
{
    \zR{承你们 老爷太太惦记着，改日再去道谢。}
}
{
    \eR{[English source not present in the checked-in Bencraft/Joly files.]}
}
{
}

\Verse{151}
{
    \zR{你们姑娘出阁，想来姑爷是不用说的了，他们的 家计如何呢？”}
}
{
    \eR{[English source not present in the checked-in Bencraft/Joly files.]}
}
{
}

\Verse{152}
{
    \zR{两个女人回道：“家计倒不怎么着，只是姑爷长的很好，为人又和 平。}
}
{
    \eR{[English source not present in the checked-in Bencraft/Joly files.]}
}
{
}

\Verse{153}
{
    \zR{我们见过好几次，看来和这里的宝二爷差不多儿，还听见说，文才也好。”}
}
{
    \eR{[English source not present in the checked-in Bencraft/Joly files.]}
}
{
}

\Verse{154}
{
    \zR{贾 母听了，喜欢道：“这么着才好，这是你们姑娘的造化。}
}
{
    \eR{[English source not present in the checked-in Bencraft/Joly files.]}
}
{
}

\Verse{155}
{
    \zR{只是咱们家的规矩还是南 方礼儿，所以新姑爷我们都没见过。}
}
{
    \eR{[English source not present in the checked-in Bencraft/Joly files.]}
}
{
}

\Verse{156}
{
    \zR{我前儿还想起我娘家的人来，最疼的就是你们 姑娘，一年三百六十天，在我跟前的日子倒有二百多天。}
}
{
    \eR{[English source not present in the checked-in Bencraft/Joly files.]}
}
{
}

\Verse{157}
{
    \zR{混的这么大了，我原想给 他说个好女婿，又为他叔叔不在家，我又不便作主。}
}
{
    \eR{[English source not present in the checked-in Bencraft/Joly files.]}
}
{
}

\Verse{158}
{
    \zR{他既有造化配了个好姑爷，我 也放心。}
}
{
    \eR{[English source not present in the checked-in Bencraft/Joly files.]}
}
{
}

\Verse{159}
{
    \zR{月里头出阁，我原想过来吃杯喜酒，不料我们家闹出这样事来，我的心就 像在热锅里熬的似的，那里能够再到你们家去？}
}
{
    \eR{[English source not present in the checked-in Bencraft/Joly files.]}
}
{
}

\Verse{160}
{
    \zR{你回去说我问好，我们这里的人都 请安问好。}
}
{
    \eR{[English source not present in the checked-in Bencraft/Joly files.]}
}
{
}

\Verse{161}
{
    \zR{你替另告诉你们姑娘，不用把我放在心上。}
}
{
    \eR{[English source not present in the checked-in Bencraft/Joly files.]}
}
{
}

\Verse{162}
{
    \zR{我是八十多岁的人了，就死 也算不得没福了。}
}
{
    \eR{[English source not present in the checked-in Bencraft/Joly files.]}
}
{
}

\Verse{163}
{
    \zR{只愿他过了门，两口儿和和顺顺的百年到老，我就心安了。”}
}
{
    \eR{[English source not present in the checked-in Bencraft/Joly files.]}
}
{
}

\Verse{164}
{
    \zR{说 着，不觉掉下泪来。}
}
{
    \eR{[English source not present in the checked-in Bencraft/Joly files.]}
}
{
}

\Verse{165}
{
    \zR{那女人道：“老太太也不必伤心。}
}
{
    \eR{[English source not present in the checked-in Bencraft/Joly files.]}
}
{
}

\Verse{166}
{
    \zR{姑娘过了门，等回了九，少 不得同着姑爷过来请老太太的安。}
}
{
    \eR{[English source not present in the checked-in Bencraft/Joly files.]}
}
{
}

\Verse{167}
{
    \zR{那时老太太见了才喜欢呢。”}
}
{
    \eR{[English source not present in the checked-in Bencraft/Joly files.]}
}
{
}

\Verse{168}
{
    \zR{贾母点头。}
}
{
    \eR{[English source not present in the checked-in Bencraft/Joly files.]}
}
{
}

\Verse{169}
{
    \zR{那女人 出去。}
}
{
    \eR{[English source not present in the checked-in Bencraft/Joly files.]}
}
{
}

\Verse{170}
{
    \zR{别人都不理论，只有宝玉听着发了一回怔。}
}
{
    \eR{[English source not present in the checked-in Bencraft/Joly files.]}
}
{
}

\Verse{171}
{
    \zR{心里想道：“为什么人家养了女 孩儿到大了必要出嫁呢？}
}
{
    \eR{[English source not present in the checked-in Bencraft/Joly files.]}
}
{
}

\Verse{172}
{
    \zR{一出了嫁就改换了一个人似的。}
}
{
    \eR{[English source not present in the checked-in Bencraft/Joly files.]}
}
{
}

\Verse{173}
{
    \zR{史妹妹这么个人，又叫他 叔叔硬压着配了人了。}
}
{
    \eR{[English source not present in the checked-in Bencraft/Joly files.]}
}
{
}

\Verse{174}
{
    \zR{他将来见了我，必是也不理我了。}
}
{
    \eR{[English source not present in the checked-in Bencraft/Joly files.]}
}
{
}

\Verse{175}
{
    \zR{我想一个人到了这个没人 理的分儿，还活着做什么！”}
}
{
    \eR{[English source not present in the checked-in Bencraft/Joly files.]}
}
{
}

\Verse{176}
{
    \zR{想到这里，又是伤心，见贾母此时才安，又不敢哭， 只得闷坐着。}
}
{
    \eR{[English source not present in the checked-in Bencraft/Joly files.]}
}
{
}

\Verse{177}
{
    \zR{一时贾政不放心，又进来瞧瞧老太太。}
}
{
    \eR{[English source not present in the checked-in Bencraft/Joly files.]}
}
{
}

\Verse{178}
{
    \zR{见是好些，便出来传了赖大，叫他将合 府里管事的家人的花名册子拿来，一齐点了一点。}
}
{
    \eR{[English source not present in the checked-in Bencraft/Joly files.]}
}
{
}

\Verse{179}
{
    \zR{除去贾赦入官的人，尚有三十馀 家，共男女二百十二名。}
}
{
    \eR{[English source not present in the checked-in Bencraft/Joly files.]}
}
{
}

\Verse{180}
{
    \zR{贾政叫现在府内当差的男人共四十一名进来，问起历年居 家用度，共有若干进来，该用若干出去。}
}
{
    \eR{[English source not present in the checked-in Bencraft/Joly files.]}
}
{
}

\Verse{181}
{
    \zR{那管总的家人将近来支用簿子呈上。}
}
{
    \eR{[English source not present in the checked-in Bencraft/Joly files.]}
}
{
}

\Verse{182}
{
    \zR{贾政 看时，所入不敷所出，又加连年宫里花用，帐上多有在外浮借的。}
}
{
    \eR{[English source not present in the checked-in Bencraft/Joly files.]}
}
{
}

\Verse{183}
{
    \zR{再查东省地租， 近年所交不及祖上一半，如今用度比祖上加了十倍。}
}
{
    \eR{[English source not present in the checked-in Bencraft/Joly files.]}
}
{
}

\Verse{184}
{
    \zR{贾政不看则已，看了急的跺脚 道：“这还了得!我打谅琏儿管事，在家自有把持，岂知好几年头里，已经‘寅年用
        了卯年’的，还是这样装好看，竟把世职俸禄当作不打紧的事，有什么不败的呢？}
}
{
    \eR{[English source not present in the checked-in Bencraft/Joly files.]}
}
{
}

\Verse{185}
{
    \zR{我如今要省俭起来，已是迟了。”}
}
{
    \eR{[English source not present in the checked-in Bencraft/Joly files.]}
}
{
}

\Verse{186}
{
    \zR{想到这里，背着手踱来踱去，竟无方法。}
}
{
    \eR{[English source not present in the checked-in Bencraft/Joly files.]}
}
{
}

\Verse{187}
{
    \zR{众人知 贾政不知理家，也是白操心着急，便说道：“老爷也不用心焦，这是家家这样的。}
}
{
    \eR{[English source not present in the checked-in Bencraft/Joly files.]}
}
{
}

\Verse{188}
{
    \zR{若是统总算起来，连王爷家还不够过的呢，不过是装着门面，过到那里是那里罢咧。}
}
{
    \eR{[English source not present in the checked-in Bencraft/Joly files.]}
}
{
}

\Verse{189}
{
    \zR{如今老爷到底得了主上的恩典，才有这点子家产，若是一并入了官，老爷就不过了 不成？”}
}
{
    \eR{[English source not present in the checked-in Bencraft/Joly files.]}
}
{
}

\Verse{190}
{
    \zR{贾政嗔道：“放屁!你们这班奴才最没良心的。}
}
{
    \eR{[English source not present in the checked-in Bencraft/Joly files.]}
}
{
}

\Verse{191}
{
    \zR{仗着主子好的时候儿，任意 开销，到弄光了，走的走跑的跑，还顾主子的死活吗？}
}
{
    \eR{[English source not present in the checked-in Bencraft/Joly files.]}
}
{
}

\Verse{192}
{
    \zR{如今你们说是没有查抄，你 们知道吗？}
}
{
    \eR{[English source not present in the checked-in Bencraft/Joly files.]}
}
{
}

\Verse{193}
{
    \zR{外头的名声，连大本儿都保不住了，还搁的住你们在外头支架子说大话， 诓人骗人？}
}
{
    \eR{[English source not present in the checked-in Bencraft/Joly files.]}
}
{
}

\Verse{194}
{
    \zR{到闹出事来，望主子身上一推就完了!如今大老爷和你珍大爷的事，说是 咱们家人鲍二吵嚷的，我看这册子上并没有什么鲍二，这是怎么说？”}
}
{
    \eR{[English source not present in the checked-in Bencraft/Joly files.]}
}
{
}

\Verse{195}
{
    \zR{众人回道： “这鲍二是不在档子上的。}
}
{
    \eR{[English source not present in the checked-in Bencraft/Joly files.]}
}
{
}

\Verse{196}
{
    \zR{先前在宁府册上。}
}
{
    \eR{[English source not present in the checked-in Bencraft/Joly files.]}
}
{
}

\Verse{197}
{
    \zR{为二爷见他老实，把他们两口子叫过 来了。}
}
{
    \eR{[English source not present in the checked-in Bencraft/Joly files.]}
}
{
}

\Verse{198}
{
    \zR{后来他女人死了，他又回宁府去。}
}
{
    \eR{[English source not present in the checked-in Bencraft/Joly files.]}
}
{
}

\Verse{199}
{
    \zR{自从老爷衙门里头有事，老太太、太太们 和爷们往陵上去了，珍大爷替理家事，带过来的，以后也就去了。}
}
{
    \eR{[English source not present in the checked-in Bencraft/Joly files.]}
}
{
}

\Verse{200}
{
    \zR{老爷几年不管家 务事，那里知道这些事呢？}
}
{
    \eR{[English source not present in the checked-in Bencraft/Joly files.]}
}
{
}

\Verse{201}
{
    \zR{老爷只打量着册子上有这个名字就只有这一个人呢，不 知道一个人手底下亲戚们也有好几个，奴才还有奴才呢。”}
}
{
    \eR{[English source not present in the checked-in Bencraft/Joly files.]}
}
{
}

\Verse{202}
{
    \zR{贾政道：“这还了得！”}
}
{
    \eR{[English source not present in the checked-in Bencraft/Joly files.]}
}
{
}

\Verse{203}
{
    \zR{想来一时不能清理，只得喝退众人。}
}
{
    \eR{[English source not present in the checked-in Bencraft/Joly files.]}
}
{
}

\Verse{204}
{
    \zR{早打了主意在心里了，且听贾赦等的官事审的 怎样再定。}
}
{
    \eR{[English source not present in the checked-in Bencraft/Joly files.]}
}
{
}

\Verse{205}
{
    \zR{一日，正在书房筹算，只见一人飞奔进来，说：“请老爷快进内廷问话。”}
}
{
    \eR{[English source not present in the checked-in Bencraft/Joly files.]}
}
{
}

\Verse{206}
{
    \zR{贾政 听了，心下着忙，只得进去。}
}
{
    \eR{[English source not present in the checked-in Bencraft/Joly files.]}
}
{
}

\Verse{207}
{
    \zR{未知吉凶，下回分解。}
}
{
    \eR{[English source not present in the checked-in Bencraft/Joly files.]}
}
{
}

\Verse{208}
{
    \zR{(本章完)}
}
{
    \eR{[English source not present in the checked-in Bencraft/Joly files.]}
}
{
}
